\documentclass[14pt,aspectratio=169,xcolor=dvipsnames]{beamer}
\usetheme{SimplePlus}
\usepackage{booktabs}
\usepackage{minted}

\title[short title]{Clase 6: Introducción a R, bash y orquestación}
\subtitle{}
\author[NA Barnafi] {Nicolás Alejandro Barnafi Wittwer}
\institute[UC|CMM] 
{
    Pontificia Universidad Católica de Chile \\
    Centro de Modelamiento Matemático
}

\titlegraphic{
    \vspace{-1.8cm}
    \begin{flushright}
      \includegraphics[height=2.5cm]{../images/logos/puc.png} 
    \end{flushright}
}
\date{}

\begin{document}
%%%%%%%%%%%%%%%%%%%%%%%%%%%%%%%%%%%%%%%%%%%%%%%%%%%%%%%
\begin{frame}
    \maketitle
\end{frame}
%%%%%%%%%%%%%%%%%%%%%%%%%%%%%%%%%%%%%%%%%%%%%%%%%%%%%%%
\begin{frame}\frametitle{Motivación}
    \begin{itemize}
        \item Distintos lenguajes tienen distintas fortalezas (Python, R)
        \item R es especializado en estadística (interpretado, sin objetos)
        \item Bash para conectar resultados (NextFlow, Snakemake)
    \end{itemize}
    \begin{flushright}
        \includegraphics[height=2cm]{../images/logos/R.png}
    \end{flushright}
\end{frame}
%%%%%%%%%%%%%%%%%%%%%%%%%%%%%%%%%%%%%%%%%%%%%%%%%%%%%%%
\section{R}
%%%%%%%%%%%%%%%%%%%%%%%%%%%%%%%%%%%%%%%%%%%%%%%%%%%%%%%
\begin{frame}[fragile]\frametitle{Ejecución de R}
    \begin{itemize}
        \item Al igual que Python, consola y programa
        \item Consola
            \begin{minted}{bash}
    $ R         # Abrir programa en bash
    > 2 + 2     # Ejecutar en consola de R
    [1] 4       # Resultado 
    > q()       # Salir de consola
            \end{minted}
        \item Programa
            \begin{minted}{bash}
    $ Rscript test.r   # Ejecutar script con '2+2'
    [1] 4              # Resultado 
            \end{minted}

    \end{itemize}
\end{frame}
%%%%%%%%%%%%%%%%%%%%%%%%%%%%%%%%%%%%%%%%%%%%%%%%%%%%%%%
\begin{frame}[fragile]\frametitle{Sintaxis: Tipos}
    \begin{itemize}
        \item String: \code{"Hola"}
        \item Números: \code{1,  3.5}
        \item Comentario: \code{"hola" \# string que saluda}
        \item Imprimir: \code{print("hola")}
        \item Asignar variable: \code{a=2, a<-2, 2->a}
        \item Variables pueden tener puntos: \code{.a=2}
        \item Sumar strings da error!
                
            \begin{flushright} \code{paste("hola","chao")}\end{flushright}
        \item Números son "numeric" por defecto (ver con \code{class})
        \item Booleanos: \code{TRUE},\code{FALSE}
    \end{itemize}
\end{frame}
%%%%%%%%%%%%%%%%%%%%%%%%%%%%%%%%%%%%%%%%%%%%%%%%%%%%%%%
\begin{frame}\frametitle{Sintaxis: Tipos II}
    \begin{itemize}
        \item Tienen nombres específicos:
            \begin{itemize}
                \item numeric \code{4,5.5,-2}
                \item integer \code{1L, -3L, 100L}  (si, con la 'L')
                \item complex \code{9+3i}
                \item character  (strings)
                \item logical   (booleanos)
            \end{itemize}
        \item Convertir números: \code{as.numeric()|as.integer()|as.complex()}
    \end{itemize}
\end{frame}
%%%%%%%%%%%%%%%%%%%%%%%%%%%%%%%%%%%%%%%%%%%%%%%%%%%%%%%
\begin{frame}\frametitle{Sintaxis II: Operaciones}
    \begin{itemize}
        \item Booleanos: \code{a\&b} (and), \code{a|b} (or), \code{!a} (not)
        \item Comparación: \code{a<b}, \code{a<=b}, \code{3==3}, \code{3!=2}
        \item Números: \code{2+2}, \code{2*3}, \code{2**3==2\^{}3}, \code{10\%\%3}
    \end{itemize}
\end{frame}
%%%%%%%%%%%%%%%%%%%%%%%%%%%%%%%%%%%%%%%%%%%%%%%%%%%%%%%
\begin{frame}[fragile]\frametitle{Sintaxis III: if/else}
    \begin{minted}{R}
    a = 2
    b = 3
    if (a>b) {
        print("a>b")
    }
    else if (a==b) {
        print("a==b")
    }
    else { # a<b
        print("a<b")
    }
    \end{minted}
\end{frame}
%%%%%%%%%%%%%%%%%%%%%%%%%%%%%%%%%%%%%%%%%%%%%%%%%%%%%%%
\begin{frame}[fragile]\frametitle{Sintaxis IV: while}
    \begin{minted}{R}
  i <- 0
  while (i < 6) {
      i <- i + 1
      if (i == 3) {
        next
      }
  print(i)
  } 
    \end{minted}
\end{frame}
%%%%%%%%%%%%%%%%%%%%%%%%%%%%%%%%%%%%%%%%%%%%%%%%%%%%%%%
\begin{frame}[fragile]\frametitle{Sintaxis V: for}
    \begin{minted}{R}
  for (x in 1:10) {
    print(x)
  } 
    \end{minted}
    \begin{minted}{R}
  fruits <- list("apple", "banana", "cherry")

  for (x in fruits) {
    print(x)
  } 
    \end{minted}
    \begin{minted}{R}
    \end{minted}
\end{frame}
%%%%%%%%%%%%%%%%%%%%%%%%%%%%%%%%%%%%%%%%%%%%%%%%%%%%%%%
\begin{frame}[fragile]\frametitle{Sintaxis VI: functions}
    \begin{minted}{R}
  my_function <- function(x, y=2) {
    return (5 * x * y)
  }

  print(my_function(3)) # y=2 por defecto
    \end{minted}
\end{frame}
%%%%%%%%%%%%%%%%%%%%%%%%%%%%%%%%%%%%%%%%%%%%%%%%%%%%%%%
\begin{frame}[fragile]\frametitle{Sintaxis VII: Estructuras de datos}
    \begin{itemize}
        \item Vector de un solo tipo: \code{c(1,2,3)},\code{1:10}
        \item Lista: \code{list(1, 2i, TRUE)}
        \item Largo: \code{length(lista)}
        \item Agregar: \code{append(lista, "hola")}
        \item Quitar: \code{lista[-1], lista[-c(1,3)]}
        \item Acceso por rango: \code{lista[2:5]}
    \end{itemize}

*OJO: Listas empiezan en 1.
\end{frame}
%%%%%%%%%%%%%%%%%%%%%%%%%%%%%%%%%%%%%%%%%%%%%%%%%%%%%%%
\begin{frame}[fragile]\frametitle{Sintaxis VIII: Data Frame}
    \begin{minted}{R}
    # Create a data frame
    Data_Frame <- data.frame (
      Training = c("Strength", "Stamina", "Other"),
      Pulse = c(100, 150, 120),
    )
    Data_Frame 

    # Todos equivalentes: 
    Data_Frame[1]
    Data_Frame[["Training"]]
    Data_Frame$Training 
    \end{minted}
\end{frame}
%%%%%%%%%%%%%%%%%%%%%%%%%%%%%%%%%%%%%%%%%%%%%%%%%%%%%%%
\begin{frame}[fragile]\frametitle{Explorando datos}
    \begin{minted}{R}
Data.Cars <- mtcars 
dim(Data.Cars)
names(Data.Cars) 
rownames(Data.Cars)
Data.Cars$cyl  # Comparar con Data.Cars["cyl"]
summary(Data.Cars)
head(Data.Cars)
max(iris["Petal.Length"]) 
which.max(iris$Petal.Length) # con [] no funciona
mean(iris$Petal.Length)
    \end{minted}
\end{frame}
%%%%%%%%%%%%%%%%%%%%%%%%%%%%%%%%%%%%%%%%%%%%%%%%%%%%%%%
\begin{frame}[fragile]\frametitle{Indexaciones y cálculos}
Si 'which' solo opera sobre vectores, cómo extraigo por nombre? \pause

    \begin{minted}{R}  
    datos = iris$"Petal.Length"
    filas = rownames(iris)
    indx_max = which.max(datos)
    filas[indx_max]
    \end{minted}
\end{frame}
%%%%%%%%%%%%%%%%%%%%%%%%%%%%%%%%%%%%%%%%%%%%%%%%%%%%%%%
\begin{frame}[fragile]\frametitle{Estadística}
    \begin{small}
\begin{minted}{R}
  model = lm(Petal.Length ~ Sepal.Length, data = iris)
  summary(model)
  b = model$coefficients[1]
  m = model$coefficients[2]
  minx = min(iris$Sepal.Length)
  maxx = max(iris$Sepal.Length)
  plot(c(minx, maxx), c(b+minx*m, b+maxx*m), 
      type="l", col="red") # (l)ine
  lines(iris$Sepal.Length, iris$Petal.Length, 
        type="p", col="blue") # (p)oint
\end{minted}
    \end{small}
    \idea{Consola...}
\end{frame}

%%%%%%%%%%%%%%%%%%%%%%%%%%%%%%%%%%%%%%%%%%%%%%%%%%%%%%%
%%%%%%%%%%%%%%%%%%%%%%%%%%%%%%%%%%%%%%%%%%%%%%%%%%%%%%%
\section{Bash}
%%%%%%%%%%%%%%%%%%%%%%%%%%%%%%%%%%%%%%%%%%%%%%%%%%%%%%%
%%%%%%%%%%%%%%%%%%%%%%%%%%%%%%%%%%%%%%%%%%%%%%%%%%%%%%%
\begin{frame}\frametitle{Navegación efectiva}
    \begin{itemize}
        \item \code{cd}, \code{ls}, etc
        \item \textbf{Usar Tab}
    \end{itemize}

    \idea{Consola...}
\end{frame}
%%%%%%%%%%%%%%%%%%%%%%%%%%%%%%%%%%%%%%%%%%%%%%%%%%%%%%%
\begin{frame}[fragile]\frametitle{Variables}
    \begin{itemize}
        \item Tienen validez solo en terminal actual
        \item Se pueden mantener en sub-terminales con \code{export}
        \item Sintaxis sin espacios
    \end{itemize}
    
    \begin{minted}{bash}
  $ A=2  # Desaparece luego de 'exit' o 'bash'
  $ export A=2 # Se mantiene dentro de nuevos procesos 'bash'
    \end{minted}
\end{frame}
%%%%%%%%%%%%%%%%%%%%%%%%%%%%%%%%%%%%%%%%%%%%%%%%%%%%%%%
\begin{frame}[fragile]\frametitle{\code{if}}
    \begin{small}
    \begin{columns}
        \begin{column}{0.45\textwidth}
            \begin{minted}{bash}
    if [ condition ]; then {
        # código por ejecutar
    }
    elif [ condition2 ]; then {
        # más código
    }
    else {
        # código final
    }
    fi
            \end{minted}
        \end{column}
        \begin{column}{0.45\textwidth}
            Acá, tenemos notación de operaciones booleanas:
            \begin{itemize}
                \item $>,\geq,<,\leq$: \code{-gt},\code{-ge},\code{-lt},\code{-le}
                \item \code{not} (\code{!}), \code{and} (\code{-a}), \code{or} (\code{-o})
                \item Ver si existe archivo/carpeta/variable: \code{-e}/\code{-d}/\code{-n}\footnote{\code{-n} ve si variable no es string vacío}
            \end{itemize}
        \end{column}
    \end{columns}
    \end{small}
\end{frame}
%%%%%%%%%%%%%%%%%%%%%%%%%%%%%%%%%%%%%%%%%%%%%%%%%%%%%%%
\begin{frame}[fragile]\frametitle{Ejemplo}
    \begin{footnotesize}
    \begin{columns}
        \begin{column}{0.7\textwidth}
            \begin{minted}{bash}
A=2
if [ ${A} -gt 1 ]; then {
    echo GT
}
else {
    echo LT
}
fi
A=""
if [ -n "${A}" ]; then { # Comillas!
    echo "SI"
}
else {
    echo NO
}
fi

            \end{minted}
        \end{column}
        \begin{column}{0.25\textwidth}
            \idea{Consola}
        \end{column}
    \end{columns}
    \end{footnotesize}
\end{frame}
%%%%%%%%%%%%%%%%%%%%%%%%%%%%%%%%%%%%%%%%%%%%%%%%%%%%%%%
\begin{frame}[fragile]\frametitle{\code{for}}
    \begin{minted}{bash}
    for n in a b c; do
       echo "$n"
    done
    \end{minted}
    Wildcards:
    \begin{itemize}
        \item \code{*.csv} $\to$ todos los archivos terminados en csv
        \item \code{*.\{csv,xls,txt\}} $\to$ todos los archivos csv, xls, o txt
    \end{itemize}

    \idea{Consola}
\end{frame}
%%%%%%%%%%%%%%%%%%%%%%%%%%%%%%%%%%%%%%%%%%%%%%%%%%%%%%%
\begin{frame}[fragile]\frametitle{Funciones}
 % Falta scripts, source, bash y los 'ifs'
    \begin{minted}{bash}
function helloWorld {
  echo "Hello World"
}
# Equivalente: 
helloWorld () {
    echo "Hello World"
}
    \end{minted}
\end{frame}
%%%%%%%%%%%%%%%%%%%%%%%%%%%%%%%%%%%%%%%%%%%%%%%%%%%%%%%
\begin{frame}[fragile]\frametitle{Funciones: argumentos}
 % Falta scripts, source, bash y los 'ifs'
    \begin{minted}{bash}
function helloWorld {
  echo "$1" # o echo "${1}"
}
    \end{minted}
    \begin{itemize}
        \item \code{\$1}, \code{\$2}, ....: parámetros en orden
        \item \code{\$0}: nombre de función
        \item \code{\$\#}: Número de argumentos   
        \item \code{\$@}: Todos los argumentos
    \end{itemize}
\idea{Consola}
\end{frame}

%%%%%%%%%%%%%%%%%%%%%%%%%%%%%%%%%%%%%%%%%%%%%%%%%%%%%%%
\begin{frame}[fragile]\frametitle{source vs bash}
    \begin{itemize}
        \item Source: Ejecutar en proceso actual
        \begin{minted}{bash}
            $ source script.sh
        \end{minted}
        \item Bash: Crear sub-proceso
        \begin{minted}{bash}
            $ sh script.sh # O 'bash script.sh'
        \end{minted}
    \end{itemize}

Consecuencias: 
    \begin{enumerate}
        \item se pueden perder variables (\code{export})
        \item  si source sale (\code{exit}), se cierra consola
    \end{enumerate}

\idea{Consola}

\end{frame}
%%%%%%%%%%%%%%%%%%%%%%%%%%%%%%%%%%%%%%%%%%%%%%%%%%%%%%%
\begin{frame}[fragile]{Pipes y redirecciones}
    \begin{footnotesize}
    \begin{minted}{bash}
  $ echo "hola"  # Imprime 'hola'
  $ # Toma el texto generado por 'echo' y lo pone en 'out.txt'
  $ echo "hola" > out.txt 
  $ # Toma el texto y lo pone al final de 'out.txt'
  $ echo "hola" >> out.txt 
  $ # Toma salida de 'cat out.txt' y lo entrega como input a grep
  $ cat 'out.txt' | grep "ho"
  $ # Tome la salida del primer comando, 1) la imprime y 2) la guarda en out.txt
  $ cat 'in.txt' | tee out.txt  # -a para 'append'
  $ # Entrega contenido de 'in.txt' a función 'echo'
  $ echo < in.txt
    \end{minted} 
    \end{footnotesize}

\idea{Consola}
\end{frame}
%%%%%%%%%%%%%%%%%%%%%%%%%%%%%%%%%%%%%%%%%%%%%%%%%%%%%%%
\begin{frame}[fragile]{Usar salidas}
    \begin{itemize}
        \item Usar el valor de la variable
            \begin{minted}{bash}
            $ echo VARIABLE="${VARIABLE}"
            \end{minted}
        \item Usar la salida de un código
            \begin{minted}{bash}
            $ A=$(python codigo.py)
            $ B=$(head "$A")
            \end{minted}
        \item Resultado de operaciones (+,-,*,/) [\code{int}]
            \begin{minted}{bash}
            $ B=3
            $ A=$((2+$B))
            \end{minted}
    \end{itemize}
\end{frame}
%%%%%%%%%%%%%%%%%%%%%%%%%%%%%%%%%%%%%%%%%%%%%%%%%%%%%%%
\begin{frame}[fragile]\frametitle{Orquestación: Ejercicio}
    Creemos: 
    \begin{enumerate}
        \item Un programa en Python que calcula el promedio de números dados
        \item Llamemos este programa desde un script en Bash
        \item Hagamos que el script en Bash abra un archivo de números y los de a Python
        \item Mostrar resultado final
    \end{enumerate}
\end{frame}
%%%%%%%%%%%%%%%%%%%%%%%%%%%%%%%%%%%%%%%%%%%%%%%%%%%%%%%
\begin{frame}\frametitle{Recap}
    \begin{itemize}
        \item En Python/R es fácil crear funciones específicas
        \item Estas funciones se pueden orquestar desde Bash
        \item Bash tiene muchas funcionalidades para manipular texto
    \end{itemize}
\end{frame}

%%%%%%%%%%%%%%%%%%%%%%%%%%%%%%%%%%%%%%%%%%%%%%%%%%%%%%%
\begin{frame}
    \maketitle
\end{frame}
%%%%%%%%%%%%%%%%%%%%%%%%%%%%%%%%%%%%%%%%%%%%%%%%%%%%%%%
\begin{frame}[noframenumbering, fragile]\frametitle{Mini ejercicios}
    \begin{tiny}
    R: 
    \begin{enumerate}
        \item Instalar R en su computador y probar consola
        \item Hacer los siguientes ejercicios: 
            \begin{center}
                https://www.w3schools.com/r/r\_exercises.asp
            \end{center}
        \item Se pueden agregar funciones no lineales de datos. Pruebe modelos no-lineales:
    \end{enumerate}

            \begin{minted}{R}
    lm(y ~ x + I(x^3) + log(x), data=datos)
            \end{minted}

Bash:
    \begin{itemize}
        \item Crear script en bash que lee la variable de ambiente 'A' y dice si es que la variable es mayor a 2.
        \item Crear script en bash que muestre el nombre y contenido de todos los archivos \code{.txt} en la carpeta.
        \item Un programa en Python que calcula el promedio de números dados
        \item Llamemos este programa desde un script en Bash
        \item Hagamos que el script en Bash abra un archivo de números y los de a Python
        \item Mostrar resultado final
    \end{itemize}

\end{tiny}
\end{frame}
\end{document}
